\begin{thema}{Δ}
Δίνεται η συνάρτηση $\displaystyle f(x) = (x^2+1)e^{-x}, \; x \in \mathbb{R}$.
\begin{erwthma}
    \item Να μελετήσετε την $f$ ως προς τη μονοτονία και να αποδείξετε ότι το σύνολο τιμών της είναι το $(0, +\infty)$.
    \item Να αποδείξετε ότι η γραφική παράσταση της $f$ έχει ακριβώς δύο σημεία καμπής, τα οποία και να βρείτε. Στη συνέχεια, να γράψετε την εξίσωση της εφαπτομένης της $C_f$ στο σημείο καμπής με τη μικρότερη τετμημένη.
    \item Να υπολογίσετε το εμβαδόν του χωρίου που περικλείεται από τη $C_f$, τους άξονες $x'x, y'y$ και την ευθεία $x=1$.
    \item Να υπολογίσετε το όριο:
    \[ L = \lim_{x\to +\infty} \frac{\ln(f(x))}{x} \]
    και στη συνέχεια να αποδείξετε ότι η εξίσωση $\displaystyle f(x) = e^{L x} + \frac{x}{4}$ έχει ακριβώς μία ρίζα στο διάστημα $(0, 1)$.
\end{erwthma}
\end{thema}
